\section{Conclusion}
\label{sec:conclusion}
% ============================================================================

The Capsule is not a decentralized database. It is not a blockchain wallet.
It is not a privacy layer bolted onto an existing system.

The Capsule is a sovereign encrypted process. It is the unit of tenancy, the
unit of privacy, the unit of collaboration, and the unit of computation in
Web5. Every principal---human, organizational, algorithmic---gets one. The
vault holds state. The identity anchors trust. The capability set governs
access. The sync topology handles replication. The compute policy determines
where code runs. The audit chain proves what happened.

Base owns state: encrypted SQLite, CRDT sync, realtime APIs, plugin runtime.
Lux owns trust: identity on I-Chain, key policy on K-Chain, threshold reveals
on T-Chain, MPC signing on M-Chain, audit anchors on the primary network.
Providers own service quality: relay throughput, storage durability, compute
availability. Users own keys and exit rights: they can decrypt, revoke,
migrate, and audit at any time, without permission from any provider.

This is what it means for the personal computer to return---not as a device
on a desk, but as a cryptographic process that follows you across every device,
every application, and every provider. The Capsule is the personal computer
of Web5.

The hard problems remain hard. Key rotation UX is unsolved for multi-device
principals. Partial sharing requires schema-level compromises. Sync metadata
leaks to traffic analysis. Private search at scale is open. Agent memory
governance needs formal verification. Cross-app composition lacks a usable
consent flow.

We have listed these problems explicitly because solving them is the work
that matters. The architecture is sound. The cryptography is proven. The
chains are specified. What remains is the engineering and UX research to
make sovereignty invisible---to make it so natural that users never think
about keys, capabilities, or sync topologies, and yet own everything.

\bibliographystyle{plain}
\begin{thebibliography}{99}

\bibitem{kleppmann2019}
M.~Kleppmann, A.~Wiggins, P.~van Hardenberg, and M.~McGranaghan,
``Local-First Software: You Own Your Data, in Spite of the Cloud,''
\textit{Onward! 2019}, ACM, 2019.

\bibitem{solid2016}
T.~Berners-Lee and D.~Sambra,
``Solid: A Platform for Decentralized Social Applications Based on Linked Data,''
MIT CSAIL Technical Report, 2016.

\bibitem{fission2021}
B.~Holmgren and P.~Chiusano,
``UCAN: User Controlled Authorization Networks,''
Fission Specification, 2021.

\bibitem{ceramic2020}
M.~Zargham, J.~Clark, and D.~Buchner,
``Ceramic: A Decentralized Data Network for Web3 Applications,''
3Box Labs, 2020.

\bibitem{dxos2023}
R.~Brinkman et al.,
``DXOS: A Framework for Local-First Collaborative Applications,''
DXOS Technical Report, 2023.

\bibitem{evervault2022}
S.~Agarwal and C.~Sheridan,
``Evervault: Encryption as Infrastructure,''
Evervault Whitepaper, 2022.

\bibitem{lux-mchain-mpc}
Z.~Kelling and V.~Seesahai,
``M-Chain: Decentralized Multi-Party Computation Custody with Quantum-Safe Threshold Signatures,''
Lux Industries, 2023.

\bibitem{lux-tfhe}
Z.~Kelling,
``Torus Threshold Fully Homomorphic Encryption: Boolean Circuits on Encrypted Data with Distributed Decryption,''
Lux Industries, 2025.

\bibitem{lux-pq-crypto}
Z.~Kelling,
``Post-Quantum Cryptographic Suite for EVM: ML-KEM, ML-DSA, and SLH-DSA as Native Precompiles,''
Lux Industries, 2024.

\bibitem{lux-fhe-mpc-hybrid}
Z.~Kelling,
``Hybrid FHE-MPC Protocols for Multi-Input Confidential Computation on Lux,''
Lux Industries, 2025.

\bibitem{lux-id-did}
Z.~Kelling,
``Lux ID: Decentralized Identifier Specification for Multi-Chain Identity,''
Lux Industries, 2024.

\bibitem{lux-tokenomics}
Z.~Kelling,
``Lux Tokenomics: Economic Design for Multi-Chain Consensus Networks,''
Lux Industries, 2023.

\bibitem{shapiro2011}
M.~Shapiro, N.~Pregui\c{c}a, C.~Baquero, and M.~Zawirski,
``Conflict-Free Replicated Data Types,''
\textit{SSS 2011}, LNCS 6976, Springer, 2011.

\bibitem{nist-mlkem}
{National Institute of Standards and Technology},
``Module-Lattice-Based Key-Encapsulation Mechanism Standard,''
FIPS 203, 2024.

\bibitem{nist-mldsa}
{National Institute of Standards and Technology},
``Module-Lattice-Based Digital Signature Standard,''
FIPS 204, 2024.

\bibitem{sqlite2022}
R.~Hipp,
``SQLite: Past, Present, and Future,''
\textit{Proc. VLDB Endow.}, 15(12), 2022.

\end{thebibliography}

